
A straight-forward reading of what has become the standard model of cosmology comes with it a strong implication that our observable universe may be just one patch of an infinite and growing multiverse
, in which the runaway expansion that shaped our deep past continues forever throughout much of space.  It is widely believed that this eternal multiverse is a generic prediction of cosmic inflation -- a free side effect arising without the need for fine-tuning of model parameters or initial conditions.  As inflation continues to gather empirical support as an explanatory account of our cosmic history, it remains an open question whether eternal inflation is truly generic among candidate models consistent with observation.

In this work, we set out to quantify the degree to which eternal inflation is a generic feature of single-field inflation models.  In other words: what is the likelihood of finding ourselves in an eternally inflating multiverse?  In addressing this question, we labor under only the most basic assumptions regarding the initial state of the inflaton and the potential function governing its dynamics.  Where quantitative analysis is no longer feasible, we attempt to paint a qualitative picture informed by mathematical considerations, and suggest steps toward developing clarity.

\paragraph{Overview}

In Section \RNum 1, we review the physics of eternal inflation and the criteria for its three distinct modes: false-vacuum, stochastic, and topological.  We discuss the implications of eternal inflation for cosmology, and take account of the established wisdom regarding the question of its prevalance among possibly predictive models of inflation.

In Section \RNum 2, we define the critera for eternal inflation to be considered {generic} (or equivalently: {\it prevalent}, {\it typical}) within the space of single-field inflation scenarios that lie within the scope of this analysis.  
We discuss the inherent non-uniqueness of the probability measure over inflaton potential functions and initial conditions, and identify some properties of the measure that we expect to enhance the statistical weight of conclusions to be drawn from it.

In Section \RNum 3, we address the prevalence of eternal inflation analytically.  
We review and expand upon several no-go theorems published in recent years, which set out criteria under which inflation is not generically eternal.  
To explore the application of these analytical constraints, we perform an in-depth analysis of hybrid inflation -- a model for which inflation is known to be non-eternal given most choices of parameters and initial conditions that are consistent with observation.

In Section \RNum 4, we take account of the gaps left in our understanding after having exhausted analytical resources, and employ numerical methods to further assess the typicality of eternal inflation.  
We simulate the dynamics of inflation across an ensemble of randomly generated potential functions and initial conditions, experimenting with several candidate measures.

In Section \RNum 5, we discuss conclusions to be drawn from this study, and suggest avenues for future work.
